\subsection{Die ursprüngliche Determinantenlinie trägt bereits primitive Windung}\label{sec:fuchswinding}
Hier liefert das Fuchssystem selbst eine weitere Herkunftsverbindung. Aus den ursprünglichen Residuen folgt
\begin{equation}
\tr A(z)=\frac{4z^3}{z^4-1},\qquad
\frac{d}{dz}\log\det\Psi(z)=\tr A(z),\qquad \Psi(0)=I.
\end{equation}
Deshalb gilt exakt
\begin{equation}
\boxed{\det\Psi(z)=1-z^4.}
\end{equation}
Die Phase $u_F=(1-z^4)/|1-z^4|$ besitzt auf dem oben konstruierten positiven Kreis um $+1$ Windung $+1$. Die tatsächliche Fuchs-Determinantenlinie stellt somit dieselbe primitive $K^1(S^1)$-Klasse dar, die P1 verlangt. Der ganzzahlige Anteil stammt aus der logarithmischen Erweiterung; aus $\det M=1$ allein könnte man ihn nicht rekonstruieren.

Die Kontur gehört zum Satz: Ein gemeinsamer Umlauf um alle vier Fuchs-Punkte hat Windung vier. Dass der ursprüngliche normale P1-Seamkreis gerade auf den einzelnen markierten Punkturkreis abgebildet wird, ist durch die bloße Klassengleichheit noch nicht bewiesen. Der nächste geometrische Herkunftsschritt ist somit eine konkrete Kontur- und Randoperatorabbildung, keine neue freie Wahl der Indexzahl.

\subsection{Die Familienoperationen wirken in der bereits vorhandenen Quelle}
Die rekonstruierte Fermionenquelle enthält $(10,1)\oplus(1,6)$ unter $\mathrm{Spin}(10)\times SU(4)$. Für $g\in SU(3)$ gilt ausdrücklich
\begin{equation}
\widehat g=\diag(1,g)\in SU(4),\qquad
\Lambda^2\widehat g\simeq g\oplus\bar g.
\end{equation}
Somit sind $M$ und $U$ wirkliche innere Symmetrien ihres Familienzweigs. Die vorhandenen sechs reellen Familienfermionen entsprechen drei komplexen Fermionen. Ihre Wirkung $M\oplus\bar M$ ist jedoch nicht $M\oplus(-M)$. Das Original-Sechserregister umfasst eine zusätzliche Sheet-Markierung. Die beiden Spin-Lifts $\widehat M$ und $-\widehat M$ haben dieselbe Vektorwirkung, aber entgegengesetzte Wirkung auf dem Familien-Spinor. Eine zusätzliche Identifikation mit dieser internen Spin-Liftwahl wäre gesondert zu beweisen. Im oben konstruierten Wegoperator entsteht das Vorzeichen bereits aus den beiden Zugängen; eine Gleichsetzung mit diesem zentralen Spin-Liftzeichen wird dafür nicht vorausgesetzt.

Der passende Anschluss an die Fuchsverbindung besteht aus Monodromiedefekten derselben Fermionenquelle. Die freie Fermionenkonstruktion bestimmt deren Antworten durch den Riemann--Hilbert-Kernel und Wick-Determinanten~\cite{gavrylenkomarshakov}. Bei der Originalkonvention $\Psi'=A\Psi$ lautet der entsprechend transponierte Kernel
\begin{equation}
\mathcal K(z,w)=\frac{[\Psi(w)^{-1}\Psi(z)]^{\mathsf T}}{z-w}.
\end{equation}
Die Defekte, ihre Sektoren und ihre Normierung sind Teile dieser Einsetzung. Eine gewöhnliche unverdrehte E8-Vakuumkorrelation wird damit nicht stillschweigend zur ursprünglichen Fuchsantwort.

Die Originalresiduen haben die Exponenten $0,1/3,2/3$. In der freien Fermionennormierung entsprechen ihnen $Q=1$ und $h=5/18$. Die gleichphasigen Vertreter $0,1/3,-1/3$ hätten dagegen $Q=0$, $h=1/9$. Gleiche Umlaufphasen erlauben keinen unbemerkten Austausch der geladenen Defektzustände. Auch der Infinity-Sektor mit Residuum $-\diag(2,1,1)$ muss erhalten bleiben. Diese genaue Markierung ist wichtiger als eine bloße Übereinstimmung der Gruppenordnungen.

Für die vollständige Randfeldwirkung ist zudem der Spinanteil getrennt zu führen. Auf dem rein internen getwisteten Schleifenraum $f(\theta+2\pi)=Mf(\theta)$ erfüllt die Halbtranslation $f(\theta)\mapsto f(\theta+\pi)$ dieselbe Sechserrelation und kommutiert mit $-i\partial_\theta$ und dessen Spektralprojektoren. Enthält die Gesamtsektion dagegen zusätzlich das NS-Spinvorzeichen, lautet ihre Randbedingung $f(\theta+2\pi)=-Mf(\theta)$ und die sechste Potenz ist $-I$. Dann wäre der rohe Halbwegdeterminant $y^6+\delta^6$. Der Original-Flavordeterminant benötigt daher die erklärte interne Transportkonvention mit getrenntem Spin-/Halbdichtefaktor oder die entsprechende andere Spinstruktur. Ein künstlich gewähltes stationäres Schleifenvakuum ersetzt nicht die tatsächlichen getwisteten Quellenmatrixelemente.

\subsection{Eine explizite gemeinsame Klasse von Randoperatoren}
Die endliche Flavorgram-Matrix besitzt außerdem eine genaue Darstellung als Dirichlet-zu-Neumann-Antwort. Für $y,\delta>0$, $y\ne\delta$, einen zyklischen Shift $S_N$ und
\begin{equation}
t=\frac{\min(y,\delta)}{\max(y,\delta)}\in(0,1),\qquad
ma=-\log t,\qquad a=\ell/N
\end{equation}
ist der Sprungoperator von $-\partial_x^2+m^2$ auf den $N$ Kreisintervallen
\begin{equation}
R_N(m)=\frac{m}{\sinh(ma)}
\bigl[2\cosh(ma)I-S_N-S_N^\dagger\bigr]
=\boxed{\frac{2m}{|y^2-\delta^2|}D_{y,\delta}^\dagger D_{y,\delta}.}
\end{equation}
Die Gleichheit folgt durch Einsetzen von $e^{-ma}=t$. Sie identifiziert eine Operatorfamilie; die Interpretation des internen $C_6$-Registers als sechs physische Kreisorte folgt daraus nicht.

Die Intervall-Determinanten und die BFK-Verklebung ergeben in dieser eindimensionalen Realisierung
\begin{equation}
2^{-N}\left(\frac{2\sinh(ma)}m\right)^N\det R_N(m)
=e^{m\ell}(1-e^{-m\ell})^2.
\end{equation}
Die lokalen Faktoren dürfen bei Variationen nicht fortgelassen werden. Die allgemeine Verklebungsformel enthält entsprechende lokale Normierungen~\cite{bfkoriginal,leebfk}. Durch Eliminieren zusätzlicher Schnittstellen reduziert dieselbe Zwölfer-Sprungmatrix exakt zur Vierer- beziehungsweise Sechsermatrix. Das ist eine gemeinsame Rechenverfeinerung desselben Kreises, keine neue Physik.

Auf dem algebraischen Originalzweig
\begin{equation}
\delta_\star=\left(\frac{794-7\sqrt{9961}}{2187}\right)^{1/6}
=0.5932772801\ldots
\end{equation}
lauten die positiven Rundläufe für $y=1,1/3,2/3$:
\begin{equation}
q_+=0.0436060052\ldots,\quad0.0314576422\ldots,\quad0.4966996535\ldots .
\end{equation}
Der elektromagnetische Rundlauf ist dagegen $q_{\rm EM}=48c_3^4e^{-2\alpha}\simeq0.0001185614$. Damit ist eine unmittelbare Identifikation aller Antworten mit derselben skalaren Massenspezialisierung ausgeschlossen. Verschiedene markierte Kanäle einer gemeinsamen Familie bleiben möglich.

Für $y<\delta$ bleibt der holomorphe Rundlauf $(\delta/y)^6>1$; die positive Darstellung benutzt seinen Kehrwert. Insbesondere gilt
\begin{equation}
\det(D^\dagger D)=\delta^{12}\bigl[1-(y/\delta)^6\bigr]^2.
\end{equation}
Der Faktor $\delta^{12}$ trägt zur Flavorvariation bei. Auch die Phase von $\det D=y^6-\delta^6$ wird durch den positiven Gramoperator allein nicht dargestellt.

\subsection{Der gemeinsame neutrale Port ist jetzt genauer bestimmt}\label{sec:neutralport}
Im Original-Mark-Greenoperator mit renormierter Diagonale null steht die dimensionslose Matrix
\begin{equation}
C_{\rm marks}=\operatorname{circ}(0,1,2,1),\qquad
G_{\rm marks}=-4c_3\log2\,C_{\rm marks}.
\end{equation}
Für den geometrischen Viererschritt $T_4$ liefert der Halbturn-gerade Projektor $E_+=(I+T_4^2)/2$ exakt
\begin{equation}
\frac14E_+C_{\rm marks}E_+=P_0,\qquad
P_0=\frac14\begin{pmatrix}1\\1\\1\\1\end{pmatrix}
\begin{pmatrix}1&1&1&1\end{pmatrix},
\qquad \det(I-qP_0)=1-q.
\end{equation}
Das benennt einen konkreten geometrischen Port. $C_{\rm marks}$ ist nicht die ursprüngliche Nahtinvolution; $E_+=P_{\Sigma,+}$ bleibt eine gesonderte Feldidentifikation.

Die Materieseite besitzt einen ladungstreuen Anschluss: Für $\mathcal R=S^+\otimes H_F$ ist
\begin{equation}
|\mathbb I_{\mathcal R}\rangle\!\rangle
=\sum_{a=1}^{48}|a\rangle\otimes|\bar a\rangle,\qquad
\|\,|\mathbb I_{\mathcal R}\rangle\!\rangle\,\|^2
=\tr_{\mathcal R}I=48
\end{equation}
der kanonische neutrale Paarungstensor. Er ist unter $V\otimes\bar V$ invariant. Dadurch können geschlossene Materieantworten in einen neutralen Port eingehen, ohne verschieden geladene Einzelteilchen zu identifizieren.

\textbf{Die fehlende gemeinsame Kopplungsidentität ist damit ausdrücklich}
\begin{equation}\label{eq:neutralselfenergy}
\boxed{\Sigma_{\rm matter}(g)
=\mathcal V^\dagger G_{\rm matter}(g)\mathcal V
=48g\,P_0,\qquad g=c_3^4e^{-2\alpha}.}
\end{equation}
Hier müssen $\mathcal V$ und $G_{\rm matter}$ aus der ursprünglichen Quelle stammen. Die neutrale Norm $48$ allein berechnet diese Selbstenergie nicht. Die Originalfassung bezeichnet $\delta_{\rm top}=48c_3^4$ selbst als Besetzungsinterpretation und definiert anschließend den Rang-eins-Transfer mit diesem Wert: \file{02_carrier_source.tex}, Zeilen 1729--1746, und \file{03_em_flavor_source.tex}, Zeilen 94--120.

Ein präziser höherer Antworttest trennt die Möglichkeiten:
\begin{align}
-\log(1-48g)&=48g+1152g^2+O(g^3),\\
-48\log(1-g)&=48g+24g^2+O(g^3).
\end{align}
Ein gemeinsamer Port und 48 unabhängige Determinanten haben dieselbe führende Spur, aber verschiedene Rückkopplungen. Die ursprüngliche zweite Defektkorrektur allein unterscheidet sie noch nicht: Mit $\beta=5/4$ beginnen beide Ausdrücke
\begin{align}
48g(1-48g)^{-\beta}&=48g+2880g^2+155520g^3+O(g^4),\\
48g(1-g)^{-48\beta}&=48g+2880g^2+87840g^3+O(g^4).
\end{align}
Der zweite Logdet-Jet beziehungsweise dritte Defekt-Jet ist deshalb der erste hier trennende gemeinsame Quellenkoeffizient. Diese Feststellung erhält die bereits berechnete Alpha-Gleichung; sie präzisiert die noch benötigte Herkunft ihrer Resummierung.

\subsection{Die tatsächliche Quelle liefert einen bestimmten neutralen Kandidaten}\label{sec:neutralfield}
Die ursprüngliche Level-eins-E8-Quelle enthält die Spinorströme $(16,4)$ mit 64 Komponenten und ihre Gegenfelder. Unter der erklärten Familienauswahl $SU(4)\to SU(3)\times U(1)$ zerfällt dieser Block in $(16,3)\oplus(16,1)$. Die erste Komponente ist der hier untersuchte 48er-Materiereadout. Seine Identifikation mit dem ursprünglichen topologischen Familienraum bleibt Teil des markierten Wörterbuchs; die Rechnung benutzt keine vierte physische Generation.

Für jeden normierten Wurzelstrom sei der symmetrische neutrale Paarvertex
\begin{equation}
B_\beta=\tfrac12\bigl(:E_\beta E_\beta^\dagger:
 +:E_\beta^\dagger E_\beta:\bigr)
=\tfrac12:(\beta\cdot h)^2:.
\end{equation}
Die zweite Gleichung folgt aus der ursprünglichen Level-eins-Gitter-OPE; $h$ sind die acht normierten Heisenbergströme. Für eine Menge $\mathcal R$ dieser Wurzeln setze
\begin{equation}
A_{\mathcal R}=\sum_{\beta\in\mathcal R}\beta\beta^{\mathsf T},\qquad
B_{\mathcal R}=\tfrac12:h^{\mathsf T}A_{\mathcal R}h:.
\end{equation}
Die exakte Rechnung an den ursprünglichen Wurzeln ergibt
\begin{align}
A_{64}&=16I_8,& B_{64}&=16T_{E_8},\\
\operatorname{spec}A_{48}&=\{12^{[5]},16^{[2]},4^{[1]}\},&
B_{48}&=12T_{D_5}+16T_{A_2}+4T_{U(1)},\\
B_{16}&=4T_{D_5}+12T_{U(1)},&
B_{48}+B_{16}&=16T_{E_8}.
\end{align}
Die Zentralzahlen der drei kommutierenden Stressanteile sind $(5,2,1)$. Das hier verwendete $U(1)$ gehört zur Familienzerlegung; es ist nicht ohne weitere Ladungsabbildung die Standardmodell-Hyperladung.

Diese Gleichungen erklären positiv, wie vollständige Quelle und 48er-Auswahl zusammengehören. Bereits die 48 geladenen Wurzelströme und ihre Gegenfelder erzeugen unter den ursprünglichen Chevalley-Klammern wieder alle E8-Wurzeln: $96\to202\to240$. Dabei kehren auch die 16 zunächst ausgelassenen Spinorwurzeln zurück. Eine Materieauslese darf diese 16 unangezeigt lassen; eine unter den ursprünglichen Feldoperationen abgeschlossene Quelle darf sie nicht löschen.

Setze $T=T_{E_8}$. Die Projektion von $B_{48}$ auf den gemeinsamen Stresszustand liefert
\begin{equation}\label{eq:nativeresidual}
\boxed{B_{48}=12T+R,\qquad R=4T_{A_2}-8T_{U(1)}.}
\end{equation}
Die Vakuumpaarungen sind
\begin{equation}
\langle T,T\rangle=4,\quad
\langle B_{48},T\rangle=48,\quad
\langle R,T\rangle=0,\quad
\boxed{\langle R,R\rangle=48.}
\end{equation}
Daher ist $\|B_{48}|0\rangle\|^2=624=576+48$.
$R$ ist ein tatsächliches neutrales Primärfeld vom konformen Gewicht zwei: Der vierte Pol seiner OPE mit $T$ verschwindet wegen $4\cdot2-8\cdot1=0$. Es ist kein Stresstensor einer neuen Theorie mit Zentralzahl null. Neutralität und Gewicht zwei allein wählen es nicht eindeutig; erst die erklärte 48er-Auswahl und die orthogonale Subtraktion des gemeinsamen Stressanteils bestimmen diese Richtung.

\picturetext{Die Antwort der 48 Materiefelder enthält einen Anteil, der das ganze Instrument gemeinsam bewegt, und eine zusätzliche Bewegung im Familienbereich. Nach Abziehen der gemeinsamen Bewegung bleibt ein konkret bestimmter Zustand. Seine Stärke ist genau 48 und seine konforme Energie genau zwei. Die beiden Zahlen stammen jetzt aus demselben Quellenoperator.}

\subsection{Ein exakter Quellenport realisiert den Alpha-Rundlauf bedingt}\label{sec:heatport}
Dieser Zustand liefert eine konkrete Realisierung der zuvor nur als Ziel formulierten Gleichung~\eqref{eq:neutralselfenergy}. Sei
\begin{equation}
|r\rangle=R_{-2}|0\rangle,\qquad
L_0|r\rangle=2|r\rangle,\qquad
\langle r|r\rangle=48,\qquad
|p_0\rangle=\tfrac12(1,1,1,1)^{\mathsf T}.
\end{equation}
Definiert man den Quellenzugang und die euklidische Quellenfortsetzung durch
\begin{equation}\label{eq:conditionalvertex}
\mathcal V=c_3^2|r\rangle\langle p_0|,\qquad
G(\alpha)=e^{-\alpha L_0},
\end{equation}
so folgt auf der vollständigen Quelle exakt
\begin{equation}\label{eq:actualheatport}
\boxed{\mathcal V^\dagger G(\alpha)\mathcal V
=48c_3^4e^{-2\alpha}P_0.}
\end{equation}
Damit liefert derselbe erklärte Port auch
$\det(I-\mathcal V^\dagger G\mathcal V)=1-48c_3^4e^{-2\alpha}$, einschließlich seiner höheren Logdet-Koeffizienten. Der Exponent zwei ist nun das Gewicht eines vorhandenen Feldzustands; die 48 ist dessen berechnete Norm. Es wird kein neuer kontinuierlicher Parameter angepasst.
Der bereits im Original verwendete Suszeptibilitätsfaktor $5/4$ wird durch diese Portidentität nicht zusätzlich hergeleitet; seine bestehende Herleitung und seine gemeinsame Normierung sind weiterhin mitzuführen.

\textbf{Status: exakte Konstruktion unter drei benannten Identifikationen, noch kein Ursprungssatz.}
Zu beweisen ist, dass die ursprüngliche Naht tatsächlich (i) den stressorthogonalen Zustand $|r\rangle$ koppelt, (ii) die Amplitude $c_3^2$ dieser Kopplung festlegt und (iii) die elektromagnetische Koordinate $\alpha$ als genau diesen konformen Fortsetzungsparameter verwendet. Die bisher geprüften Originaldefinitionen von $P_{\rm prim}$, $P_{\rm sing}$ und $P_\Theta$ geben diese Abbildung nicht an. Insbesondere bedeutet die ursprüngliche primitive Windungsbedingung nicht automatisch die Projektion auf ein Virasoro-Primärfeld.

Der Port ist eine Abbildung in einen bestimmten Quellenzustand. Er ist nicht automatisch dasselbe wie eine beliebige lokale Kopplung an das Feld $R(z)$. Dessen Dreipunktantwort lautet
\begin{equation}
\langle R(z_1)R(z_2)R(z_3)\rangle
=\frac{-384}{z_{12}^2z_{23}^2z_{31}^2}.
\end{equation}
Diese nichtverschwindende Antwort muss bei lokalen Mehrfacheinsetzungen mitgeführt werden. Sie schließt außerdem eine vakuumerhaltende Quellenautomorphie aus, die $R$ schlicht auf $-R$ abbildet. Eine solche Nahtparität kann diesen Anteil nicht ohne weiteres löschen. Der vollständige Quellprozess bleibt reicher als seine eine geometrische Portanzeige.

\subsection{Flavor und neutrale Quelle besitzen bereits eine gemeinsame berechnete Antwort}\label{sec:jointneutralward}
Der neue Kandidat lässt sich unmittelbar am ursprünglichen Flavor-Hintergrund prüfen. Dafür wird die normierte Riemann--Hilbert-Defektantwort derselben drei komplexen Familienfermionen verwendet, mit den ursprünglichen Residuen und dem vorgeschriebenen ausgehenden Sektor. In dieser Darstellung liefern die Fermionen-Wardidentitäten
\begin{equation}
\langle T_{\rm Familie}(z)\rangle_F=\tfrac12\tr A(z)^2,
\qquad
\langle T_{U(1)}(z)\rangle_F=\tfrac16\bigl(\tr A(z)\bigr)^2.
\end{equation}
Die symmetrische Stressdefinition entfernt den Ableitungsterm; der verbundene $U(1)$-Stromzweipunktteil hat weiterhin den Koeffizienten drei. Da $T_{\rm Familie}=T_{A_2}+T_{U(1)}$, erhält man durch Einsetzen der ursprünglichen Verbindung
\begin{equation}\label{eq:jointRresponse}
\boxed{\langle R(z)\rangle_F
=2\bigl[\tr A(z)^2-(\tr A(z))^2\bigr]
=\frac{-20z^6+(52/9)z^2}{(z^4-1)^2}.}
\end{equation}
Das ist eine konkrete Antwort des neutralen Kandidaten auf die bereits vorhandenen Flavor-Defekte. Sie wurde unabhängig symbolisch gegengeprüft. Da $R$ ein Primärfeld vom Gewicht zwei ist, ist diese Antwort als quadratisches Differential zu transformieren; eine zusätzliche Schwarzian-Korrektur entsteht für $R$ nicht.

Die lokalen Doppelpolkoeffizienten an jeder der vier Punkturen betragen $-8/9$. Am unendlichen Ende gilt $z^2\langle R(z)\rangle_F\to-20$. Die negativen Werte sind kein Positivitätsverstoß: $R$ ist eine Differenz von Stressantworten und kein positiver Energieoperator. Die Originalresiduen $(0,1/3,2/3)$ liefern lokal
\begin{equation}
h_{A_2}=1/9,\qquad h_{U(1)}=1/6,\qquad
4h_{A_2}-8h_{U(1)}=-8/9.
\end{equation}
Der lokale Rest bei Unendlich besitzt die Exponenten $(-2,-1,-1)$; im ausgehenden Bra-Sektor entsprechen sie der Ladung $+4$. Er kompensiert die vier Defekte mit Ladung $+1$. Triviale Monodromie am unendlichen Ende bedeutet somit nicht Vakuumladung. Genau diese ganzzahlige Quelleninformation ginge bei einer Beschreibung ausschließlich durch die Monodromiematrizen verloren.

Auch der erste lokale Normtest ist jetzt berechenbar. Für einen normierten Highest-Weight-Defektzustand $|\sigma\rangle$ mit den obigen Gewichten folgt aus der Virasoro-Relation
\begin{equation}\label{eq:defectRnorm}
\|R_{-2}|\sigma\rangle\|^2
=16(4h_{A_2}+1)+64(4h_{U(1)}+1/2)
=\boxed{880/9},
\end{equation}
während die Vakuumnorm 48 beträgt. Eine bloße Mittelwertsubtraktion ändert den Zustand $R_{-2}|\sigma\rangle$ nicht. Das ist ein lokaler Sektortest, keine Behauptung, dass die gesamte Vierdefektgeometrie aus einem einzigen Highest-Weight-Zustand besteht.

\key{\textbf{Was der gemeinsame Befund entscheidet:}
Der Flavortransport und ein exakter neutraler Alpha-Port sind nun innerhalb der vorhandenen Quellenstruktur konkret darstellbar. Ihre gemeinsame Antwort ist jedoch zustands- und markierungsabhängig und bereits ausdrücklich berechnet. Die gesuchte Nahtabbildung muss deshalb erklären, welcher Quellenkanal ausgewertet wird und wie seine Normierung mit den Flavor-Einsetzungen zusammenhängt. Verschiedene Einsetzungen desselben Quellenfunktionals sind zulässig; ihre gemeinsamen Antworten dürfen nicht durch getrennte Vakuumannahmen ersetzt werden.}

\subsection{Die vollständige neutrale Antwort im ursprünglichen Flavorprozess}\label{sec:fullflavorRR}
Die gemeinsame Einpunktantwort bestimmt bereits einen Teil der Verträglichkeit.
Der vollständige verbundene Zweipunktkern liefert nun die tatsächliche Antwort
zwischen zwei verschiedenen Einsetzungsorten. Entscheidend ist die unveränderte
Feldkonvention des ursprünglichen Fuchs- und Riemann--Hilbert-Systems:
\[
 \Psi'(z)=A(z)\Psi(z),\qquad j(z)=\tr A(z),\qquad
 \mathcal K(z,w)=[\Psi(w)^{-1}\Psi(z)]^{\mathsf T}/(z-w).
\]
Im normierten quasifreien Riemann--Hilbert-Defektmatrixelement folgt aus
symmetrischer Fermion-Punktaufspaltung~\cite[Abschnitte 3.3 und 4.1]{gavrylenkomarshakov}
\begin{align}
 \langle T_{\rm Familie}(z)T_{\rm Familie}(w)\rangle_{F,c}
 &=\frac{3}{2(z-w)^4}
   +\frac{\tr\!\left(A(z)A(w)\right)}{(z-w)^2},\\
 \langle T_{U(1)}(z)T_{U(1)}(w)\rangle_{F,c}
 &=\frac{1}{2(z-w)^4}+\frac{j(z)j(w)}{3(z-w)^2}.
\end{align}
Die $A_2$- und $U(1)$-Faktorisierung ergibt für $R=4T_{\rm Familie}-12T_{U(1)}$
\begin{equation}\label{eq:fullflavorRR}
\boxed{\begin{aligned}
 \langle R(z)R(w)\rangle_{F,c}
 &=\frac{48}{(z-w)^4}\\[-2pt]
 &\quad+\frac{16}{(z-w)^2}
 \left[\tr\!\left(A(z)A(w)\right)+j(z)j(w)\right].
\end{aligned}}
\end{equation}
Der zusätzliche Term ist für die ursprüngliche Vierdefektverbindung vollständig rational:
\begin{equation}
 \langle R(z)R(w)\rangle_{F,c}-\frac{48}{(z-w)^4}
 =\frac{32(99z^3w^3+4z^2+5zw+4w^2)}
 {9(z-w)^2(z^4-1)(w^4-1)}.
\end{equation}
Das Indexzeichen $c$ bedeutet, dass das Produkt der Einpunktwerte bereits
abgezogen wurde. Der Zusatzterm ist also keine bloße Verschiebung des Mittelwerts.

Der lokale Doppelpolkoeffizient lautet exakt
\begin{equation}
 C_2(z)=16\left[\tr A(z)^2+j(z)^2\right]
 =\frac{32z^2(99z^4+13)}{9(z^4-1)^2}.
\end{equation}
Er stimmt mit der unabhängigen OPE-Identität
$C_2=64\langle T_{\rm Familie}\rangle_F-8\langle R\rangle_F$ überein.
Jede der vier Punkturen liefert $224/9$; dadurch wird auch die lokale Norm
$48+2(224/9)=880/9$ aus~\eqref{eq:defectRnorm} wiedergewonnen.
Bei $w=0{,}43$ und $z=0{,}17$ ist der verbundene Kern numerisch
$10572{,}47509$, gegenüber $10503{,}83390$ im Vakuum.

Die Reihenfolge im Kernel ist wesentlich. Die Primärquelle verwendet
$\varphi'=\varphi A_{\rm rechts}$. Für unsere linke Fuchs-Gleichung gilt
$\varphi=\Psi^{\mathsf T}$ und $A_{\rm rechts}=A^{\mathsf T}$.
In der neutralen Wick-Antwort kürzen sich die Fundamentallösungen deshalb
vollständig; es bleibt $\tr(A(z)A(w))$.
Die zunächst untersuchte vertauschte Matrixreihenfolge hätte stattdessen
$\tr(A(z)P(z,w)A(w)P(w,z))$ geliefert. Sie gehört zu einer anderen lokalen
Feldrahmung und wird hier ausdrücklich nicht als ursprüngliche RHP-Antwort
übernommen. Die korrigierte Antwort ist rational und einwertig, wie es die
Invarianz von $R$ unter den Familienmonodromien verlangt.

Das Programm \file{tfpt_explorer/source_flavor_correlator.py} berechnet
jetzt diese eindeutige Antwort aus der Originalverbindung. Die
Reihenfolgenprüfung gehört zur Regression; lokale Einpunkt- und
OPE-Ergebnisse bleiben unverändert. Die ausgewiesene Reichweite ist das
ursprüngliche Defektmatrixelement. Eine positive globale Dichtematrix oder
physische Photonpolarisation wird daraus nicht stillschweigend abgeleitet.

\key{\textbf{Die neue gemeinsame Bedingung:}
Eine Nahtabbildung im Flavor-Hintergrund muss den gesamten Kern
\eqref{eq:fullflavorRR} in der richtigen ursprünglichen Feldkonvention abbilden.
Das bloße Ersetzen durch die Vakuumzahl 48 ist unzureichend.
Ein vollständiger $E_+$-Port mit Referenzdeterminante kann trotzdem zulässig
sein: Die gerade lokale Korrektur allein entscheidet diese globale,
normalisierte Abbildung noch nicht.}

\picturetext{Der Familienweg bewegt sich innerhalb desselben Instruments.
Deshalb verändert er auch, wie dieses Instrument auf zwei neutrale Anfragen
gemeinsam reagiert. Diese Veränderung ist jetzt ausgerechnet. Sie gehört zur
gemeinsamen Mechanik, die eine vollständige Lösung erhalten muss.}

\paragraph{Der ursprüngliche Randport ist bereits berechenbar.}
Die Originalverifikationen v484 und v485 enthalten den Viermarken-Kern
und seinen logarithmischen Referenzabzug. Ihre Matrix lautet
\[
 C_{\rm mark}=\begin{pmatrix}
 0&1&2&1\\1&0&1&2\\2&1&0&1\\1&2&1&0
 \end{pmatrix},\qquad
 K_{\rm mark}=-c_3\log2\,C_{\rm mark}.
\]
Hier bezeichnet $K_{\rm mark}$ bereits das Orbitmittel des ursprünglichen
Mark-Greenkerns. Mit der zyklischen Markverschiebung $Q_4$,
$E_{\rm mark,+}=(I+Q_4^2)/2$ und $p_0=(1,1,1,1)^{\mathsf T}/2$ folgt exakt
\[
 E_{\rm mark,+}C_{\rm mark}E_{\rm mark,+}=4P_0,\quad
 P_0=p_0p_0^\dagger,\quad
 p_0^\dagger K_{\rm mark}p_0=-4c_3\log2.
\]
Der gerade Markraum hat Dimension zwei; die komprimierte Antwort hat Rang eins.
In der Originalkonvention $u=\varepsilon c_3\log2$ sind
\begin{equation}
 \det(I-uC_{\rm mark})=(1-4u)(1+2u)^2,\qquad
 \det(I-uE_{\rm mark,+}C_{\rm mark}E_{\rm mark,+})=1-4u.
\end{equation}
Die schon vorhandene Alpha-Vorgabe $q=48c_3^4e^{-2\alpha}$ würde dadurch
exakt getroffen, wenn der ursprüngliche Kontaktparameter die Identifikation
\begin{equation}\label{eq:originalportmatching}
 \boxed{\varepsilon(\alpha)=\frac{12c_3^3}{\log2}e^{-2\alpha}}
\end{equation}
besitzt. Dies ist eine aus beiden vorhandenen Antworten berechnete
\emph{Matchingbedingung}; sie wird hier nicht als unabhängige Herleitung der
Kontaktkopplung ausgegeben.

Die jetzt präzisierte Anschlussfrage lautet daher: Welche ursprüngliche
Variation $\partial_\alpha D_\xi$ beziehungsweise Rand-zu-$R$-Einsetzung
erzeugt~\eqref{eq:originalportmatching} und denselben Referenzabzug im
Flavor-Hintergrund? v485 subtrahiert die logarithmische Singularität von
$|D|^{-1}$. Das definiert noch keine endliche Diagonale des neuen
$R$-$R$-Kerns vierter Polordnung an den Fuchsdefekten.
Ebenso ist der Mark-Halbdrehungsprojektor nicht allein durch die Bezeichnung
mit dem ursprünglichen Calderón- oder Nahtprojektor identifiziert.
Der fehlende Gegenstand ist diese konkrete Einsetzungs- und
Referenzabbildung, nicht ein generell unbekannter ursprünglicher Randkernel.
Die Formeln stammen aus \file{verification/v484_seam_contact_unit.py},
\file{verification/v485_contact_diagonal_closed.py} und dem
Alpha-Transfer in \file{_archive/tfpt-45/source_extracts/03_em_flavor_source.tex}.

